\section{星表位置与自行的修正方法}
过去，已有研究提出通过比较不同星表与高精度参考星表之间的恒星位置差异，
在局部天区尺度上刻画并修正星表系统误差。
其中，Chesley 等人基于 2MASS 星表的工作奠定了区域去偏方法的基础，
而随后 Farnocchia 等人将该框架扩展至同时包含位置与自行修正。
在 Gaia 高精度星表发布之后，Eggl 等人进一步以 Gaia 数据作为统一参考，
系统性地更新并完善了历史星表的去偏方案。

本工作总体遵循 Eggl 等人提出的 Gaia 时代星表去偏思想，
以 Gaia DR3 作为统一的天球参考框架，
通过比较历史星表与 Gaia DR3 中恒星源的位置与自行，
在局部天区尺度上估计星表的系统性偏差，
并将所得修正量应用于历史小天体光学观测的校正中。
与既有工作的主要区别在于，
本研究在工程实现层面针对大规模星表数据进行了系统优化，
并在统一处理流程下同时涵盖多套历史星表及多个 Gaia 数据版本。

由于星表系统误差在天球上通常表现出明显的空间相关性，
其统计建模需要在有限的角尺度上进行。
为此，全天球被划分为若干等面积天区，
并假定在每一个天区内星表的系统偏差可以近似视为常量。
在本分析中，采用 HEALPix 等面积像素化方案将全天球划分为 $49{,}152$ 个天区，
对应的单个天区面积约为 $0.8~\mathrm{deg}^2$。

为了定量评估不同历史星表相对于 Gaia DR3 的系统偏差，
不同星表之间应具有一致的比较基准，
在本研究中，所有恒星位置首先被统一转换至 J2000.0 历元，
对于上世纪甚至更老的星表，还需要考虑参考系的统一。
之后在每一个天区内对目标星表与 Gaia DR3 之间的恒星源进行空间交叉匹配。

% 整体的匹配逻辑与eggl相似，在本研究中细化了参数的选择与匹配判据。
% 具体来说，对于星表而言，对GAIA DR3 中的恒星源进行星表过滤可以显著减低暗星的干扰，
% 从而提升匹配的可靠性与精度。
% 对星表中的每一颗待匹配源而言，匹配通过在参考位置附近设定搜索半径来完成，

并在存在多个候选匹配时施加严格的判据以消除歧义，
仅保留角距离显著优于其他候选的唯一匹配。
此外，强制施加一一对应约束，
以避免单个 Gaia 恒星源与多个目标星表源发生关联，
从而降低密集星场中伪匹配和多对一匹配
对系统误差统计结果的污染。

在完成可靠的恒星交叉匹配后，
对每个天区内匹配恒星的位置与自行差异进行统计分析。
所有差异均在 J2000.0 历元下计算，
分别对应赤经与赤纬方向。
对于每一个天区和每一个目标星表，
分别估计其在 J2000.0 历元下的位置修正
$(\Delta\alpha_{2000}, \Delta\delta_{2000})$
以及对应的自行修正
$(\Delta\mu_{\alpha}, \Delta\mu_{\delta})$。
为降低异常值、残余误匹配以及局部离群点的影响，
上述修正量均采用稳健统计量（中位数）进行估计。
所得结果刻画了目标星表相对于 Gaia DR3 的局部系统性偏差结构。

对于一条在历元 $t$ 给出的小天体光学观测，
其由参考星表系统误差引入的修正量可表示为
\begin{equation}
\Delta\alpha(t) = \Delta\alpha_{2000} + \Delta\mu_{\alpha}\,(t - 2000.0),
\end{equation}
\begin{equation}
\Delta\delta(t) = \Delta\delta_{2000} + \Delta\mu_{\delta}\,(t - 2000.0),
\end{equation}
其中赤经方向的修正项已包含球面几何中的 $\cos\delta$ 因子。
将上述修正从原始观测中扣除后，
所有历史光学观测均可视为已统一至 Gaia DR3 参考框架下，
从而能够在一致的误差模型中用于轨道确定与残差分析。

\newpage